\documentclass[10pt]{article}
\usepackage[margin=0.9in]{geometry}
\usepackage{amsmath,microtype}
\title{Guide: All Irrational Mechanical Itineraries Are Excluded}
\author{Collatz proof project}
\date{September 5, 2026}
\begin{document}
\maketitle
\section*{What Lean now proves}
No natural Collatz orbit can have the parity pattern
\[
\lfloor(t+1)\alpha+\rho\rfloor-\lfloor t\alpha+\rho\rfloor
\]
when $\alpha$ is irrational. This holds for every real intercept $\rho$,
and even if the pattern starts after any finite number of orbit steps.
Earlier results treated particular slopes and an infinite family. The new
proof covers every irrational slope in this mechanical-word definition.

\section*{Why the argument works}
A parity block is a finite string recording whether successive orbit values
are even or odd. An unbounded Collatz orbit must produce sufficiently many
different blocks: identical long blocks force early starting values to be
equal, which would trap the orbit in a repeating cycle.

A mechanical pattern has too few blocks. For length $L$, slide its starting
intercept through one unit interval. There are only $L$ floor thresholds
that can change the resulting prefix counts. An ordered rank from zero to
$L$ distinguishes the possible blocks, so there are at most $L+1$.
Combining these bounds forces the orbit to be bounded.

A bounded natural orbit repeats a state. Suppose the repeating section has
$q$ steps and contains $p$ odd steps. Repeating it $k$ times contributes
exactly $kp$ odd steps. The mechanical formula puts this within one of
$kq\alpha$. This is possible for all $k$ only if $q\alpha=p$, making the
slope rational. That contradicts the assumed irrationality.

\section*{What this changes, and what remains}
The formal bridge between the complexity proof and the repository's
slope-and-intercept definition is now complete. The argument is a
formalization of classical ideas; no claim of mathematical novelty is made.
The technical paper gives the exact declarations and literature context.

This does not prove Collatz. A possible counterexample need not have a
mechanical parity pattern. For rational mechanical slopes, this proof only
forces boundedness; it does not prove that the eventual cycle is the usual
one. The central remaining problem is a restriction covering every possible
counterexample, or a general descent proof.

\section*{Reproduce the verification}
From \texttt{lean/}, run
\texttt{lake build Collatz.MechanicalComplexity}.
From the repository root, with Lean on the executable path, run
\texttt{python3 analysis/audit\_theorems.py --build}.
The audit checks selected declaration dependencies, not a separate kernel
implementation. Both paper and guide have editable LaTeX sources.
\end{document}
